\documentclass[11pt]{article}
\usepackage{amsmath,amssymb}
\usepackage[margin=1in]{geometry}
\usepackage{booktabs}
\usepackage{hyperref}
\usepackage{xcolor}
\hypersetup{colorlinks=true, linkcolor=blue!50!black, urlcolor=blue!50!black,
            citecolor=blue!50!black}

\newcommand{\Psib}{\Psi}
\newcommand{\Var}{\mathcal{V}}
\newcommand{\Mem}{\mathcal{M}}
\newcommand{\Ex}{\mathcal{E}}
\newcommand{\code}[1]{\texttt{#1}}
\DeclareMathOperator{\Sol}{Sol}
\DeclareMathOperator{\ld}{ld}
\DeclareMathOperator{\ord}{ord}

\newtheorem{definition}{Definition}
\newtheorem{example}{Example}
\newtheorem{proposition}{Proposition}
\newtheorem{lemma}{Lemma}
\newtheorem{remark}{Remark}
\newtheorem{criterion}{Criterion}
\newtheorem{principle}{Principle}

\title{Prolongation, projection, and the quantifier discipline:\\
       an algorithm for the ansatz-plus-PDE solution space\\
       and the resolution of the $v = v_1x+v_2y+v_3z$ puzzle}
\author{}
\date{\today}

\begin{document}
\maketitle

\begin{abstract}
\noindent
A companion note to \emph{A New Method of Solving PDEs}. Given an ansatz $A(c)$
(a differentially triangular family with constant parameters $c$) and a PDE
$P$, we survey the algorithm families available for computing the solution
space of the combined problem, then develop a proposed three-phase algorithm:
prolong $A \cup \{P\}$ to differential closure while tracking the inequations
actually used (Phase~I), recurse Thomas-style on the tracked inequation loci
(Phase~II), and project each resulting system onto the constants under the
quantifier discipline \emph{independent variables universal, dependent
variables and constants existential} (Phase~III). The intellectual center of
the note is a puzzle: applied naively, the universal coefficient-zeroing of
Phase~III would read the ansatz-defining equation $v = v_1x+v_2y+v_3z$ as
forcing $v_1=v_2=v_3=0$. We resolve the puzzle by making the quantifier
discipline precise. Universal quantification is implemented by
\emph{coefficient collection} and is sound only over symbols that are
algebraically free of one another --- which the fully reduced remainder's
variables are, and a defining equation's leader is not; existential
quantification is implemented by \emph{Gr\"obner elimination} and is sound only
\emph{after} differential closure (prolongation to passivity). The naive
Phase~III interleaves the two operations in an order in which they do not
commute, and we exhibit small examples where it over-constrains (the
$v$-equation) and where it under-constrains (returning $\mathbb{C}$ where the
correct existence locus is $\{c=1\}$). Corrected, the algorithm computes the
per-cell \emph{existence} locus --- the joint-variety projection
$\Var \cup \{\text{partial strata}\}$ --- with the \emph{membership} locus
$\Var$ selectable per cell by flipping the quantifier on the parametric
derivatives, equivalently by the dimension-preservation criterion.
\end{abstract}

\tableofcontents

%=====================================================================
\section{Problem statement: two loci, not one}
\label{sec:problem}

Let $A(c) = \{a_1,\dots,a_k\}$ be the \emph{ansatz}: a differentially
triangular set of differential polynomials in dependent variables
$\Psib, v, \dots$ over independent variables $x,y,z$ (derivations
$\partial_x,\partial_y,\partial_z$), whose coefficients involve constant
parameters $c \in \mathbb{C}^n$. Let $P$ be the PDE. The running concrete case
is the hydrogen ansatz of the paper \cite{NewMethod}:
\begin{equation}
\label{eq:hydrogen-ansatz}
\begin{gathered}
(a_0+a_1v)\,\Psib'' + (b_0+b_1v)\,\Psib' + (c_0+c_1v)\,\Psib = 0,\\
v - (v_1x+v_2y+v_3z+v_4r) = 0, \qquad r^2 - (x^2+y^2+z^2) = 0,\\
\Psib_x - \Psib'v_x = 0,\quad \Psib_y - \Psib'v_y = 0,\quad \Psib_z - \Psib'v_z = 0,
\end{gathered}
\end{equation}
with $P$ the Schr\"odinger equation
$-\tfrac12(\Psib_{xx}+\Psib_{yy}+\Psib_{zz}) - \tfrac1r\Psib = E\Psib$ and
eleven constants $c = (E, v_1,\dots,v_4, a_0,a_1,b_0,b_1,c_0,c_1)$.

``The solution space of ansatz $+$ PDE'' admits two readings, and the entire
subject of this note turns on keeping them distinct
\cite{membership-note}:
\begin{align}
\Mem &\;=\; \{\, c : P \in [A(c)] \,\}
  && \text{(\emph{membership locus}: every ansatz member solves $P$),}
  \label{eq:M}\\
\Ex &\;=\; \pi_c\bigl(\Sol(A)\cap\Sol(P)\bigr)\setminus\{\text{trivial}\}
  && \text{(\emph{existence locus}: some nonzero member solves $P$).}
  \label{eq:E}
\end{align}
Here $[A(c)]$ is the differential ideal generated by the ansatz at $c$, and
``trivial'' excludes the identically-zero solution, which for the
linear-homogeneous class solves everything at every $c$ and must be quotiented
away before either locus is informative
\cite[\S2bis, ``Two saturations'']{bounded-prolongation-note}.

$\Mem$ is a \emph{universal} condition over the ansatz family; $\Ex$ is
\emph{existential}. Always $\Mem \subseteq \Ex$, and the difference
\[
  \Ex \setminus \Mem \;=\; \{\text{\emph{partial-solution strata}}\}
\]
consists of parameter values where a \emph{proper nontrivial subfamily} of the
ansatz solves $P$ while the family as a whole does not
\cite[Def.~1]{membership-note}. The paper's Algorithm~1 (reduce $P$ modulo
$A$, project the remainder onto the constants) computes $\Mem$; decomposing the
combined system $A \cup \{P\}$ and projecting computes $\Ex$. The two-line
example
\[
A:\ \Psib''-c\,\Psib, \qquad P:\ \Psib'-\Psib
\]
separates them: $\Mem = \varnothing$ (no $c$ makes the whole two-dimensional
family $K_1e^{\omega x}+K_2e^{-\omega x}$ obey $\Psib'=\Psib$), yet
$\Ex = \{c=1\}$ (the subfamily $Ke^{x}$ solves at $c=1$)
\cite[Ex.~1]{membership-note}.

The paper's own motivation for treating $A$ and $P$ \emph{on an equal
footing} --- one combined decomposition rather than reduce-then-project ---
is precisely that the partial strata are invisible to the membership
formulation by construction, and that some of them carry physically meaningful
solutions: in the hydrogen differential Thomas run, the bad-locus cell at
$E=-1/8$ is the hydrogen $2s$ state $(1-r/2)e^{-r/2}$, whose constants
provably fail the projection system of Algorithm~1
\cite{existence-route-report}. Any algorithm that adjoins $P$ to $A$ must
therefore be judged against \emph{both} targets: which locus does it compute,
and can it be tuned to compute the other?

%=====================================================================
\section{Survey: the algorithm families}
\label{sec:survey}

Four routes to the ansatz-plus-PDE solution space are on the table. The
conceptual comparison --- what question each answers and what triggers a case
split --- is developed at length in \cite{method-comparison}; we summarize,
adding only what this note needs. Prior-art verdicts and full citations are in
the litsearch report \cite{litsearch}.

\subsection{Reduce-then-project (the paper's Algorithm 1)}

Ritt's full reduction of $P$ modulo $A$ over the field $\mathbb{Q}(c)$
produces $h\cdot P - r \in [A]$ with $r$ fully reduced and $h$ a power product
of initials and separants of $A$ (the ``denominator''). The projection step
treats each surviving non-constant symbol as an independent indeterminate,
collects
\begin{equation}
\label{eq:collection}
r \;=\; \sum_{\alpha=1}^{N} m_\alpha\,p_\alpha(c),
\end{equation}
$m_\alpha$ distinct power products in the non-constants, and returns
$\Var = V(p_1,\dots,p_N)$ as minimal primes (commutative GTZ). For hydrogen
the collection runs over exactly $x,y,z,r,\Psib,\Psib'$
\cite[\S{}Projection Step]{NewMethod}, yields 28 constant equations, and the
five minimal primes reproduced by \code{joca.sage}.

\emph{Question answered:} the membership locus, $\Var = \Mem$ (on the locus
$h \ne 0$; Theorem~1's completeness hypothesis). \emph{Splits:} none in the
good case --- one reduction, one GTZ call; the general Algorithm~2 recurses
only on the \emph{bad locus} where the hypotheses of the completeness theorem
fail. \emph{Cost/faithfulness:} the constants never enter the differential
engine (they live in the coefficient field), which is the practical lever ---
and the price is that the bad-locus strata (e.g.\ $a_0=a_1=0$, carrying the
$1s$ and $2s$ states as first-order-ODE solutions) are outside the guarantee.

The closest published relative is Chaharbashloo--Basiri--Rahmany--Zarrinkamar
\cite{chaharbashloo}: substitute an \emph{explicit} physics ansatz
$\psi = g(r)e^{y(r)}$ into the wave equation, equate powers of $r$ by hand,
and solve the resulting constant system by a commutative Gr\"obner basis. That
is exactly a reduce-then-project computation --- and, as \S\ref{sec:puzzle}
will make precise, the reason their raw power-equating is sound is that an
explicit ansatz has no defining equations left: substitution has already
eliminated every dependent leader.

\subsection{Differential Thomas decomposition}

B\"achler--Gerdt--Lange-Hegermann--Robertz \cite{BGLR2012} (implementation for
PDE systems: the \textsc{Maple} package TDDS \cite{TDDS}; project Sage port in
progress, \code{joca-thomas-native-dt.sage} / \code{thomas\_cells.py}).
Decomposes any system into finitely many \emph{disjoint} simple systems
(cells): each differentially triangular, \emph{passive} (Janet-involutive:
every non-multiplicative prolongation reduces to zero), with attached
inequations (initials, separants, squarefreeness witnesses).

\emph{Question answered:} the complete, disjoint stratification of
$\Sol(A\cup\{P\})$ --- the existential object; the constant-space shadow of the
cells is $\Ex$ together with its case structure. With the constants modeled as
$\dot c_i = 0$, the constant-only inequations become parameter-degeneracy
conditions and the decomposition is automatically comprehensive over the
constants. \emph{Splits:} wherever the \emph{description} degenerates ---
every initial, separant, or discriminant of every equation processed during
completion, whether or not its vanishing changes anything about existence.
\emph{Status here:} the open-maple run of the hydrogen combined system
completes with 29 cells; the bad-locus cells 26/27 carry the $E=-1/2$ and
$E=-1/8$ strata \cite{existence-route-report}.

\subsection{Parametric Rosenfeld--Gr\"obner}

Fakouri--Rahmany--Basiri \cite{fakouri}; the non-parametric engine is
Boulier's BLAD/BMI implementing Rosenfeld--Gr\"obner \cite{blop}. Regular
differential chains represent $[C]:H_C^\infty$; a comprehensive-Gr\"obner-style
bookkeeping splits parameter space wherever an initial or separant's vanishing
is ambiguous mid-elimination.

\emph{Question answered:} per parameter cell, a regular-chain description of
the (saturated) solutions --- existential in character, but of the
\emph{saturation} $[C]:H_C^\infty$, so strata inside initial/separant loci are
systematically divided out rather than kept. \emph{Splits:} parameter
conditions arising during elimination. \emph{Status here:} the project port
(\code{\char`~/parametric-rg}) reproduces the paper's examples; on hydrogen
the parameter branching is cheap but the differential reduction hits the BLAD
coefficient-swell wall \cite{method-comparison}. Saturation-unfaithfulness is
not hypothetical: reducing the PDE against the RG chain's bare equations
re-admits the spurious prime $(a_0,a_1)$ while absorbing (losing) the two
genuine strata inside it \cite[\code{joca-rg.sage} discussion]{README}.

\subsection{Bounded prolongation to the involutive order, then pure
commutative algebra}
\label{sec:bounded-prolongation}

The route developed in \cite{bounded-prolongation-note}: prolong
$A\cup\{P\}$ to an order $N$ at which the system is involutive, treat every
jet of order $\le N$ as an independent algebraic indeterminate, form the
\emph{un-saturated} polynomial ideal
$J = ((A\cup\{P\})^{[N]}) \subseteq K[\text{jets}_{\le N}, x,y,z,r, c]$, and
finish with any commutative-algebra algorithm. The differential engine is
touched exactly once. The semantic core is published (Lange-Hegermann's
Theorem~4.1 \cite{LH2014}: truncated solution sets of an involutive system
equal the solution sets of its algebraized truncation; the construction
itself appears inside Golubitsky--Kondratieva--Ovchinnikov's treatment of the
generalized Ritt problem \cite{GKO2009}, there with the bad-locus--deleting
saturation $:H^\infty_C$ that this route deliberately omits). Comprehensive
analogues on the algebraic side: comprehensive involutive systems
\cite{gerdt-hashemi} and comprehensive triangular decomposition
\cite{CTD2007,chen-morenomaza}.

\emph{Question answered:} the existence locus,
$V\bigl((J:\Psib^\infty)\cap K[c]\bigr) = \overline{\Ex}$ --- the
$:\Psib^\infty$ quotient (delete only the zero section) is mandatory, and is
a different operation from the $H_A$-saturation the route avoids
\cite[\S2bis]{bounded-prolongation-note}. \emph{Splits:} optional --- a
constructible (algebraic-Thomas) split on initials/separants/discriminants
recovers the disjoint strata; the single-ideal route returns the Zariski
closure, a superset that never loses a component. \emph{Two corrections from
the July 2026 implementation attempt} \cite{existence-route-report}, both
inherited by the algorithm of \S\ref{sec:algorithm}:
\begin{itemize}
\item The involutive order of the \emph{ansatz alone} ($N_0(A) = 2$ for
hydrogen, from the single-$v$ structure) is \textbf{not} the involutive order
of the \emph{combined} system: the PDE--ODE critical pair (shared leader
$\Psib''$) feeds conditions $D^k\rho$ (derivatives of the Ritt remainder) into
the ideal only at order $2+k$. No a-priori $N_0(A\cup\{P\})$ is known;
\emph{prolong-until-stable} --- watch the elimination ideal
$(J^{(N)}:\Psib^\infty)\cap K[c]$ ascend and stop at stabilization (Hilbert
basis theorem) --- is the algorithm, not a safety check.
\item The direct implementation over all eleven constants as ring variables
hits a comprehensive-elimination cliff (no completion even over
$\mathrm{GF}(32003)$ at $N=2$); per-cell staging over the already-computed
Thomas cells, or pointwise Nullstellensatz probes, are the practical
fallbacks.
\end{itemize}

\subsection{Summary table}

\begin{center}
\small
\begin{tabular}{p{2.9cm}p{3.4cm}p{3.6cm}p{3.6cm}}
\toprule
 & question answered & what triggers a split & principal limitation \\
\midrule
Algorithm 1 (reduce{+}project) & membership $\Mem$ on $\{h\ne0\}$ &
  only bad-locus recursion (Alg.~2) & misses partial strata and the
  $h=0$ strata by construction \\
differential Thomas & full disjoint stratification of
  $\Sol(A\cup\{P\})$; shadow $=\Ex$ & every initial / separant /
  discriminant met during completion & splitting cost; constants ride
  through the differential engine \\
parametric RG & per-cell regular chains of the \emph{saturated} ideal &
  parameter ambiguity mid-elimination & saturation deletes in-locus
  strata; coefficient swell \\
bounded prolongation & $\overline{\Ex}$ (closure), strata optional &
  optional constructible split & no a-priori $N_0(A\cup\{P\})$;
  parametric elimination cliff \\
\bottomrule
\end{tabular}
\end{center}

%=====================================================================
\section{The proposed algorithm}
\label{sec:algorithm}

We now state the three-phase proposal this note is asked to develop, first as
proposed, then --- after the analysis of \S\S\ref{sec:quantifier}--%
\ref{sec:puzzle} --- in corrected form (\S\ref{sec:corrected}).

\subsection{Phase I: prolong, reduce, track}

Prolong $A\cup\{P\}$ and auto-reduce until differential closure --- every
integrability condition (non-multiplicative prolongation, equivalently
$\Delta$-polynomial) reduces to zero. By the first correction of
\S\ref{sec:bounded-prolongation}, this must be run as a completion loop
(prolong $\to$ reduce $\to$ adjoin any nonzero residue, terminating by
Ritt--Raudenbush), not as a single prolongation to a precomputed bound: for
the combined system the constants-cutting conditions are the reduced
derivatives $D^k\rho$ of the Ritt remainder, and no a-priori bound on $k$ is
available.

During every reduction step, record which reductor was used and how:
\emph{algebraic} steps divide by the reductor's \emph{initial};
\emph{differential} steps (reduction against a proper prolongation
$\theta a$) divide by its \emph{separant}. In the algebraized picture the two
cases unify: the initial of every proper prolongation $\theta a$ \emph{is}
the separant of $a$, so the tracked set is simply ``the initials of the
(possibly prolonged) equations actually used.'' The product of the tracked
divisors is exactly Ritt's denominator: the reduction identity is
\begin{equation}
\label{eq:reduction-identity}
h \cdot p \;\equiv\; r \pmod{(T^{[N]})},
\qquad h = \prod_i I_i^{d_i}\,\prod_j S_j^{e_j},
\end{equation}
a polynomial identity, valid whether or not $h$ vanishes; only the
\emph{conclusions} drawn from $r$ (e.g.\ ``$r=0$ hence $p$ is a consequence'')
require $h \ne 0$.

\subsection{Phase II: recurse on the tracked inequations}

For each tracked inequation $h_i$, split:
\[
\Sol(\Sigma) \;=\; \Sol(\Sigma \cup \{h_i \ne 0\})
   \;\sqcup\; \Sol(\Sigma \cup \{h_i = 0\}),
\]
keep $h_i \ne 0$ on the main branch, and recurse Phase~I on the branch with
$h_i$ adjoined as an equation. Termination: adjoining $h_i$ strictly enlarges
the (radical differential) ideal, and strictly ascending chains of radical
differential ideals terminate (Ritt--Raudenbush); when $h_i \in K[c]$ the
split strictly descends the parameter dimension, the same descent that
terminates Algorithm~2's bad-locus recursion and CTD's discriminant recursion.
The output is a finite tree of systems $(E_i = 0,\ NE_i \ne 0)$, each
differentially closed (passive) relative to its inequations.

This is the Thomas/RG splitting discipline with lazy bookkeeping: Thomas also
splits only on initials, separants, and squarefreeness conditions of equations
it actually processes --- not on arbitrary potential zero-divisors --- so
Phases I--II are best understood as a differential Thomas decomposition whose
split set is generated \emph{on demand} by the reduction trace. Two honest
caveats. First, ``used'' must include the separants that license the
\emph{parametric-data freeness} of the final cell (each solved-form equation's
initial must be a unit for the Riquier existence theorem to apply); these are
the separants of the final triangular set, which the completion loop does use,
but an implementation that tracked only mid-reduction divisors and not
final-form initials would emit cells on which Phase~III's existence semantics
silently fail. Second, disjointness of the output requires carrying the
complementary equations honestly ($h_i = 0$ branches), exactly as Thomas does;
dropping them recovers Algorithm~1's ``main branch plus bad locus'' shape.

\subsection{Phase III as proposed --- and why it cannot stand as stated}
\label{sec:phase3-naive}

The proposal: in each equation of each leaf system, \emph{group like terms by
the independent variables and their algebraic extensions} ($x,y,z,r$ with
$r^2 = x^2+y^2+z^2$); because the equation must hold for all real values of
the independents, set every grouped coefficient to zero; each coefficient
involves constants and dependent variables; since a dependent variable need
only take \emph{some} value, \emph{eliminate} the dependent variables from the
coefficient system by Gr\"obner elimination and return the constant varieties
admitting at least one solution in the dependents.

The organizing intuition --- \textbf{independents universally quantified,
dependents and constants existentially quantified} --- is semantically
correct, and \S\ref{sec:quantifier} will keep it. But the operational
recipe ``collect over the independents first, then eliminate the dependents
inside the coefficients'' fails on both sides, and the failures are not
edge cases:

\begin{example}[over-constraint: the defining equation]
\label{ex:v-naive}
Apply the recipe to the ansatz equation $v - v_1x - v_2y - v_3z = 0$.
Grouping by $x,y,z$ gives four coefficient rows: $v$ (the constant term),
$-v_1$, $-v_2$, $-v_3$. Three rows contain no dependent variable at all, so
the existential step has nothing to witness and the recipe forces
$v_1 = v_2 = v_3 = 0$ --- destroying a perfectly good \emph{defining}
equation of the ansatz. (The $v$-row is ``witnessed'' by $v=0$, compounding
the damage.)
\end{example}

\begin{example}[under-constraint: the frozen witness]
\label{ex:frozen}
Apply the recipe to the combined toy system $A: \Psib''-c\Psib$,
$P: \Psib'-\Psib$ \emph{without} the Phase-I closure (or, equivalently, apply
it directly to Algorithm~1's reduced remainder $r = \Psib'-\Psib$). No
independent variable appears, so there is a single coefficient row
$\Psib'-\Psib$, containing dependents only; the existential step asks whether
\emph{some} values $(\psi_0,\psi_1)$ satisfy $\psi_1-\psi_0 = 0$ --- always
yes. The recipe returns all of $\mathbb{C}$, but $\Ex = \{c=1\}$ and
$\Mem = \varnothing$. The recipe is not computing the existence locus: it has
quantified over \emph{frozen} values of $\Psib,\Psib'$, constant in $x$,
whereas a solution's jets vary with $x$ --- the constraint that kills
$c \ne 1$ is the \emph{derivative} of the imposed relation,
$D(\Psib'-\Psib) = \Psib''-\Psib' \to c\Psib - \Psib = (c-1)\Psib$, which
only differential closure supplies.
\end{example}

Example~\ref{ex:v-naive} over-constrains ($\{v_1=v_2=v_3=0\}$ where the true
answer is ``no condition''); Example~\ref{ex:frozen} under-constrains
($\mathbb{C}$ where the true answer is $\{1\}$). So the naive Phase~III
computes a locus that is in general \emph{incomparable} to both $\Mem$ and
$\Ex$: it is the set of constants for which some $x$-constant assignment of
the dependent symbols annihilates the equations identically in the
independents --- a hybrid with no differential meaning. The next two sections
diagnose why, and what the correct discipline is.

%=====================================================================
\section{The quantifier principle, made precise}
\label{sec:quantifier}

\subsection{The four classes of symbols}

Fix a leaf system of Phase~II: a differentially triangular set $T$ with
inequations $NE$, under a ranking $<$ on the jets. The non-constant symbols
fall into four classes, and \emph{each class admits exactly one sound
projection operation}:

\begin{definition}[symbol classes]
\label{def:classes}
Relative to $(T, <)$:
\begin{enumerate}
\item \textbf{Constants} $c$: the answer space. Never quantified away; the
  output lives in $K[c]$.
\item \textbf{Independent variables} $x,y,z$ (and their algebraic extensions,
  here $r$ with $r^2 = x^2+y^2+z^2$): universally quantified --- an equation
  must hold at every point of some open set.
\item \textbf{Principal jets}: the leaders of $T$ and all their derivatives.
  Each is \emph{determined} by its equation (solved form; the initial is a
  unit by $NE$): its quantifier is irrelevant because it has a unique witness.
\item \textbf{Parametric jets}: the remaining jets --- the free Taylor data of
  the family cut out by $T$ (finite in number for a finite-type system).
  \emph{This class is the selector}: quantify it universally and the
  projection computes membership; existentially, existence.
\end{enumerate}
\end{definition}

For the hydrogen ansatz under the paper's ranking
$\Psib'' > \Psib' > \Psib > v > r > \text{constants} > x,y,z$: the principal
jets are $\Psib''$, $v$, $r$ (leaders), the spatial derivatives $\Psib_x$,
$\Psib_y$, $\Psib_z$ (leaders of the chain rules), and all their derivatives;
the parametric jets are exactly $\Psib, \Psib'$. It is no accident that the
paper's projection collects like terms in precisely
$x,y,z,r,\Psib,\Psib'$ --- classes 2 and 4 --- and in nothing else.

\subsection{The two implementations of $\forall$, and the one of $\exists$}

\begin{principle}[operations, not just quantifiers]
\label{prin:ops}
Universal quantification is implemented by \emph{coefficient collection};
existential quantification by \emph{elimination}; determined symbols by
\emph{substitution} (triangular rewriting). Collection is sound only over
symbols algebraically independent of one another and of everything retained;
elimination over the jets is sound only after differential closure.
\end{principle}

We justify the three clauses in turn.

\paragraph{Collection ($\forall$).}
An equation $e \in K[c][x,y,z,r][\text{parametric jets}]$, containing no
principal jet, vanishes for all real values of the independents and all
admissible values of the parametric data iff every collected coefficient
$p_\alpha(c)$ in \eqref{eq:collection} vanishes. Two ingredients:
\begin{itemize}
\item \emph{Syntactic:} distinct power products are linearly independent over
  $K[c]$, so $e \equiv 0$ as a polynomial iff all $p_\alpha = 0$. For the
  algebraic extension $r$ this requires $e$ to be reduced modulo
  $r^2-(x^2+y^2+z^2)$ first ($r$-degree $\le 1$); the reduced monomials are
  linearly independent as functions on the quadric. This is one of several
  places where ``fully reduced'' is load-bearing.
\item \emph{Semantic (density):} vanishing \emph{for all real} values is
  equivalent to identical vanishing because $\mathbb{R}^3$ is Zariski-dense in
  $\mathbb{C}^3$; vanishing for all admissible parametric data is equivalent
  because, on a cell with its inequations, the parametric Taylor data is
  freely assignable at every admissible base point (Riquier existence --- for
  the hydrogen class, elementary second-order-ODE existence), so the realized
  values of $(x,y,z,r,\Psib,\Psib')$ fill a Zariski-dense subset of the
  quadric-times-jet space.
\end{itemize}
Note what collection is \emph{not}: it is not elimination. One does not
eliminate a universally quantified variable from an ideal --- elimination
computes the projection of a variety, which is the $\exists$ semantics.
Collection is the linear-algebra operation dual to it, and this asymmetry is
why the constants can stay in the coefficient field while the universal
variables stay symbolic.

\paragraph{Elimination ($\exists$).}
For the existence reading, the object is the prolonged, algebraized,
un-saturated ideal of the leaf system,
$J = (T^{[N]}) \subseteq K[\text{jets}_{\le N}, x,y,z,r, c]$, at closure
order $N$. Lange-Hegermann's Theorem~4.1 \cite{LH2014} (equivalently Lemma~3
of \cite{bounded-prolongation-note}) says that \emph{for an involutive
system}, every point of $V(J)$ respecting the inequations is the truncation
of a genuine formal solution. Therefore
\begin{equation}
\label{eq:elimination}
J_c \;=\; \Bigl(J : \bigl(\Psib\cdot\textstyle\prod_{h\in NE}h\bigr)^{\!\infty}\Bigr)
      \cap K[c]
\end{equation}
cuts out the Zariski closure of the cell's contribution to $\Ex$: eliminate
\emph{all} the jets \emph{and the base point} $x,y,z,r$ (a solution needs only
\emph{some} point of analyticity --- demanding all points would wrongly
exclude, e.g., the Coulomb singularity $r=0$), after quotienting out the zero
section and localizing at the cell's inequations. The universal
``for all points near the base point'' is carried entirely by the
prolongation generators $\theta a \in T^{[N]}$: prolongation is the
Taylor-basis implementation of $\forall$-over-independents, dual to the
monomial-basis implementation by collection. Example~\ref{ex:frozen} is
exactly what happens when this clause is ignored: without the generators
$D^k(\Psib'-\Psib)$, the eliminated ideal knows nothing of $(c-1)\Psib$ and
the projection overshoots.

\paragraph{Substitution (determined).}
A principal jet is rewritten by its solved-form equation --- this is what
Ritt reduction \emph{is}. Its equation is a definition, not a constraint row;
its unique witness is supplied by the triangular structure. Feeding a
defining equation into either projection operation is the category error of
Example~\ref{ex:v-naive}.

\subsection{The two sound pipelines}

Assembling Principle~\ref{prin:ops}, there are exactly two sound orders of
operations --- one per quantifier assignment on the parametric jets --- and
the naive Phase~III is sound in neither reading because it interleaves them
\cite[Finding~3: ``exactly one $\forall$ algorithm and one $\exists$
algorithm'']{existence-route-report}:

\begin{center}
\small
\begin{tabular}{lll}
\toprule
 & membership $\Mem$ (per cell) & existence $\Ex$ (per cell) \\
\midrule
1. differential closure & full reduction of $P$ by $T$ &
  prolongation of $T\cup\{P\}$ to passivity \\
2. principal jets & substituted away by the reduction &
  eliminated (they are determined) \\
3. parametric jets & \textbf{collected} (universal) &
  \textbf{eliminated} (existential, after 1!) \\
4. independents $x,y,z,r$ & collected (universal) &
  eliminated ($\exists$ base point; $\forall$ carried by 1) \\
5. constants & land in $K[c]$: the system $\{p_\alpha(c)\}$ &
  land in $K[c]$: the ideal $J_c$ of \eqref{eq:elimination} \\
\bottomrule
\end{tabular}
\end{center}

In the membership column no coefficient containing a dependent symbol is ever
``set to zero'': collection continues \emph{through} the parametric jets until
the coefficients are pure constants. In the existence column nothing is ever
collected: every non-constant is eliminated, and the universal content lives
in the prolongation generators. The trap is precisely the hybrid --- collect
over the independents, stop, and existentially witness the jet-carrying
coefficients --- because collection over $x$ silently reclassifies every
dependent symbol as a constant-in-$x$ coefficient, which dependents are not
(they are functions of $x$; that is what ``dependent'' means).

\begin{remark}[relation to comprehensive Gr\"obner systems and QE]
The $\forall\exists$ block structure over an algebraically closed field is a
quantifier-elimination problem, and the tools align: the $\exists$ block is
constructible-image computation (elimination ideals, per parameter cell ---
the comprehensive Gr\"obner system / Gr\"obner cover technology
\cite{weispfenning,montes-wibmer}), while the $\forall$ block, because it is
universal over a \emph{polynomial identity}, degenerates to coefficient
collection --- linear, requiring no cell analysis and no real QE. This is why
the whole pipeline stays inside Gr\"obner-land. The only real-vs-complex
subtlety sits in the $\exists$ block: elimination certifies a \emph{complex}
witness; realness of solutions is a separate question the paper handles at
the level of the final varieties.
\end{remark}

%=====================================================================
\section{Resolution of the $v = v_1x+v_2y+v_3z$ puzzle}
\label{sec:puzzle}

\subsection{The distinguishing criterion}

\begin{criterion}
\label{crit:main}
An equation may participate in independent-variable coefficient-zeroing if
and only if it contains \emph{no principal jet} --- equivalently, iff it is
fully reduced with respect to the cell's triangular set --- and the collection
must then run over the independents \emph{and} the parametric jets together,
terminating in coefficients built from constants alone. An ansatz-defining
equation never qualifies: its leader is a principal jet \emph{by
definition} --- the equation is the solved form that \emph{defines} that jet
--- so it enters the projection only as a rewrite rule (membership reading)
or as an elimination generator (existence reading), never as a constraint row
to be collected.
\end{criterion}

Applied to the puzzle: in $v - v_1x - v_2y - v_3z = 0$ the symbol $v$ is a
dependent variable --- itself a function of $x,y,z$ --- and under any
admissible ranking it is the equation's leader (it is the only jet present).
The equation is therefore a definition of $v$, and Criterion~\ref{crit:main}
excludes it from collection in both readings:

\begin{itemize}
\item \emph{Membership reading (Algorithm 1):} $v$ is principal, so full
  reduction substitutes $v \mapsto v_1x+v_2y+v_3z\,(+\,v_4r)$ everywhere in
  the remainder before the projection ever looks at it. The defining equation
  is consumed as a rewrite rule; $v$ never appears among the collection
  variables --- witness the paper's collection set $x,y,z,r,\Psib,\Psib'$,
  from which $v$ is conspicuously absent while the \emph{parametric} jets
  $\Psib,\Psib'$ are present.
\item \emph{Existence reading (bounded prolongation):} the equation and its
  prolongations $v_x - v_1$, $v_y - v_2$, $v_z - v_3$, $v_{xx}, \dots = 0$
  enter the ideal $J$, and the $v$-jets are \emph{eliminated}. The elimination
  ideal $\bigl(v - v_1x - v_2y - v_3z,\ v_x-v_1,\ \dots\bigr) \cap
  K[v_1,v_2,v_3]$ is zero: no condition on the constants, as it must be.
\end{itemize}

Either endpoint is safe. The naive Phase~III fails because it stops the
universal collection at the independents, leaving ``coefficients'' that still
contain dependent symbols, and then treats those coefficients as objects that
must vanish (rows without dependents --- forcing $v_1=v_2=v_3=0$) or be
witnessed by frozen values (rows with dependents --- Example~\ref{ex:frozen}).
Both failure modes are the same error seen from two sides: \emph{collection
over $x,y,z$ presupposes that everything left in a coefficient is constant in
$x,y,z$, and dependent variables are not.}

\subsection{What distinguishes the defining equation from the reduced
remainder}

Both $v - v_1x-v_2y-v_3z$ and a reduced remainder $r$ are elements of the same
differential ring; both may contain jets; the remainder $r = \Psib'-\Psib$ of
the toy example even has a jet \emph{leader} ($\Psib'$). The distinction is
not ``contains a jet'' but the \emph{status of its jets relative to the
triangular set doing the projecting}:
\begin{itemize}
\item the defining equation contains its own leader --- a jet that is
  principal \emph{for the ansatz} --- so relative to $A$ it is a definition;
\item the fully reduced remainder contains, by the definition of full
  reduction, no leader of $A$ and no derivative of one: every jet in it is
  \emph{parametric} for $A$, i.e.\ free family data over which the membership
  question universally quantifies.
\end{itemize}
The ranking is what decides which jet of each equation is its leader, hence
which jets are principal, hence which symbols the collection may treat as
free. Change the ranking and the same polynomial can change roles ---
which is not a defect but the design: the ranking of the paper
(\S{}Ranking of \cite{NewMethod}) is engineered exactly so that $\Psib''$,
the spatial derivatives, $v$, and $r$ are the defined symbols and
$\Psib,\Psib'$ the free ones.

The ``equal footing'' aspiration is then realized without contradiction:
adjoining $P$ to $A$ and decomposing means the remainder itself \emph{becomes
a defining equation} of the combined system --- in the toy example,
$\Psib'-\Psib$ acquires the leader $\Psib'$, turning $\Psib'$ from parametric
(for $A$) to principal (for $A\cup\{P\}$), and the completion then derives
$(c-1)\Psib$, whose \emph{initial} $(c-1)$ carries the constant condition as a
Thomas split, not as a coefficient-zeroing. On an equal footing, constants
conditions materialize as \emph{initials of derived equations} and as
\emph{elimination ideals of cells} --- never by zeroing jet-carrying
coefficients. Coefficient-zeroing reappears only when, a cell having been
fixed, one asks the membership question \emph{of that cell} --- at which point
its defining equations act again as rewrite rules and
Criterion~\ref{crit:main} applies verbatim.

\begin{remark}[why the physics tradition never met the puzzle]
The ansatz of \cite{chaharbashloo} (and of the quasi-exact-solvability
literature it implements) is \emph{explicit}: $\psi = g(r)e^{y(r)}$ is
substituted into the wave equation before any algebra happens, so every
dependent symbol is eliminated by construction and raw power-equating is
Criterion~\ref{crit:main}-compliant without anyone having to say so. The
puzzle is a phenomenon of \emph{implicit} (differentially triangular)
ans\"atze, where definitions coexist with constraints in one system and
must be told apart --- by their leaders.
\end{remark}

\subsection{The criterion tested on the partial-strata examples}
\label{sec:tests}

\paragraph{Example A ($A: \Psib''-c\Psib$, $P: \Psib'-\Psib$; one derivation).}
\emph{Membership pipeline:} $P$ is already fully reduced w.r.t.\ $A$ (its
leader $\Psib'$ ranks below $\Psib''$, and no derivative of $\Psib''$
appears); its jets $\Psib,\Psib'$ are parametric for $A$; collect over them:
coefficients $+1, -1$; $\Mem = V(1,-1) = \varnothing$. \checkmark\ (matches
\cite[Ex.~1]{membership-note}).
\emph{Existence pipeline:} close: $(A\cup\{P\})^{[2]} =
\{\Psib''-c\Psib,\ \Psib'-\Psib,\ \Psib''-\Psib'\}$; auto-reduce:
$\Psib''-\Psib' \to c\Psib - \Psib' \to (c-1)\Psib$, tracking no division
(all initials are $1$ until the derived equation appears); the derived
equation's initial $(c-1)$ splits Phase-II-style: on $c \ne 1$ the cell
forces $\Psib = 0$ (trivial; discarded by the $:\Psib^\infty$ quotient), on
$c = 1$ the cell is $\{\Psib' = \Psib\}$ with parametric jet $\Psib$;
eliminating jets with $\Psib \ne 0$ leaves the ideal $(c-1)$, i.e.\
$\Ex = \{1\}$. \checkmark\ The dimension criterion agrees: the $c=1$ cell has
one parametric jet where $A$ has two --- a proper subfamily, hence a partial
stratum, hence in $\Ex \setminus \Mem$.

\paragraph{Example B ($A: \Psib_{yy}-c\Psib$, $P: \Psib_x-\Psib$; two
derivations).}
Here $A\cup\{P\}$ is already autoreduced \emph{and coherent}: the
$\Delta$-polynomial of the leader pair $(\Psib_x, \Psib_{yy})$ reduces to
zero, so Phase~I terminates with no derived equation and Phase~II never
splits --- one cell, equations $\{\Psib_x-\Psib,\ \Psib_{yy}-c\Psib\}$,
parametric jets $\{\Psib, \Psib_y\}$.
\emph{Membership:} reduce $P$ by $A$: remainder $P$ itself; collect over the
parametric data of $A$ (which is infinite --- $A$ alone is not finite type;
the collection over $\Psib_x, \Psib$ suffices): coefficients $+1,-1$;
$\Mem = \varnothing$. \checkmark
\emph{Existence:} eliminate the jets from the closed cell with
$\Psib \ne 0$: the prolongations
($\Psib_{xy} = \Psib_y$, $\Psib_{xyy} = \Psib_{yy} = c\Psib$ from both
routes, \dots) are consistent for every $c$; the elimination ideal in $K[c]$
is $(0)$; $\Ex = \mathbb{C}$ --- every $c$ is a partial stratum, the stratum
being the subfamily $e^x(K_1e^{\omega y}+K_2e^{-\omega y})$.
\checkmark\ (matches \cite[Ex.~2]{membership-note}). The dimension criterion
again separates the readings: the cell's parametric set $\{\Psib,\Psib_y\}$
is a collapse from $A$'s infinite parametric set (two arbitrary functions
$F(x), G(x)$ down to two constants), so no cell is dimension-preserving and
$\Mem = \varnothing$ --- even though the decomposition never split and the
adjoined equation ``reduced to nothing new,'' confirming that membership
cannot be read off the splitting combinatorics, only off the quantifier
discipline (or dimensions).

\paragraph{The puzzle equation itself.} Worked in
\S\ref{sec:puzzle}: no condition on $v_1,v_2,v_3$ in either pipeline;
$v_1=v_2=v_3=0$ arises only from the unsound hybrid. \checkmark

%=====================================================================
\section{The corrected algorithm, and which locus it computes}
\label{sec:corrected}

\subsection{Statement}

\begin{quote}
\textbf{Algorithm (prolongation--projection, corrected).}
\emph{Input:} ansatz $A(c)$, PDE $P$, ranking $<$.
\emph{Output:} the existence locus $\Ex$ stratified into cells, with each
cell's membership status decided; hence both loci.
\begin{enumerate}
\item[\textbf{I.}] \textbf{Close and track.} Run the completion loop on
  $A\cup\{P\}$ (prolong $\to$ auto-reduce $\to$ adjoin nonzero residues) to
  passivity, tracking every initial/separant used as a divisor
  \eqref{eq:reduction-identity}. No a-priori order bound is assumed
  (prolong-until-stable; the ansatz-only bound $N_0(A)$ is \emph{not} valid
  for the combined system).
\item[\textbf{II.}] \textbf{Split and recurse.} For each tracked divisor
  $h_i$: keep $h_i \ne 0$ on the current branch; spawn a branch with
  $h_i = 0$ adjoined and recurse from I. Collect the leaves
  $(T_j, NE_j)$ --- finitely many, each passive.
\item[\textbf{III$_\exists$.}] \textbf{Project each leaf onto the constants
  (existence).} Algebraize ($\text{jet} \mapsto$ indeterminate) at the leaf's
  closure order; compute
  $J_c^{(j)} = \bigl(J_j : (\Psib\prod_{h\in NE_j}h)^\infty\bigr)\cap K[c]$,
  eliminating all jets and the base point. Then
  $\overline{\Ex} = \bigcup_j V\bigl(J_c^{(j)}\bigr)$, each leaf contributing
  its stratum (bad-locus strata included --- nothing was saturated by $H_A$
  globally).
\item[\textbf{III$_\forall$.}] \textbf{Decide membership per leaf
  (optional selector).} On each leaf, either (a) fully reduce $P$ by $T_j$
  and collect over independents \emph{and} the leaf's parametric jets,
  testing whether the leaf's constant relations already imply every collected
  coefficient (radical membership, as implemented in
  \code{joca-rg-combined.sage}); or equivalently (b) test dimension
  preservation: the leaf's parametric-derivative set equals the ansatz's.
  Leaves passing are membership cells; their union recovers $\Mem = \Var$;
  the rest are partial strata.
\end{enumerate}
\end{quote}

The changes from the proposal of \S\ref{sec:algorithm} are exactly two, both
forced by \S\ref{sec:quantifier}: Phase~I is a completion loop rather than a
one-shot prolongation to $N_0(A)$; and Phase~III's two quantifier moves are
\emph{unmixed} --- elimination of everything (III$_\exists$) or collection
through everything non-constant (III$_\forall$), never collection-then%
-elimination on the same equation.

\subsection{Which locus --- and the selector}

Phase III$_\exists$ computes, per leaf, the Zariski closure of that leaf's
contribution to the existence locus; the union over leaves is
$\overline{\Ex} = \overline{\Var \cup \{\text{partial strata}\}}$
\cite[\S2bis(b), superset remark]{bounded-prolongation-note}. It is a
\emph{new route to the joint-variety projection}, not to a third locus: by
the two-pipelines analysis there is no third quantifier assignment to
compute. What is genuinely selectable is the per-cell refinement:

\begin{itemize}
\item \textbf{Existence per cell} (III$_\exists$): keeps the bad-locus
  strata (hydrogen $E=-1/2$, $E=-1/8$ cells) that Algorithm~1's
  $h\ne0$ hypothesis fences off --- the whole point of the equal-footing
  treatment.
\item \textbf{Membership per cell} (III$_\forall$): the
  \emph{dimension-preservation criterion} of
  \cite[\S{}The correct selection criterion]{membership-note} --- a cell
  belongs to $\Var$ iff the PDE removed nothing from the family, i.e.\ iff
  the cell's parametric set is the ansatz's. Lange-Hegermann's differential
  counting polynomial \cite{LH2014} makes this effective by inspection of
  the cell (count the parametric derivatives); the radical-membership test of
  \code{joca-rg-combined.sage} is the implemented equivalent. Note the test
  must be the \emph{radical-membership} form, not ``re-reduce $P$ by the
  cell's chain and collect'': reduction against a chain's bare equations does
  not saturate, and the hydrogen \code{joca-rg.sage} run shows the resulting
  projection re-admitting the spurious prime $(a_0,a_1)$ while absorbing the
  genuine strata inside it \cite{README}.
\end{itemize}

So the corrected algorithm is a \emph{selectable hybrid}: III$_\exists$ alone
gives $\Ex$; III$_\exists$ filtered by III$_\forall$ gives $\Var$; reporting
both, with each partial stratum labeled by its cell, gives strictly more
information than either endpoint algorithm --- which is the honest way to
state its value over Algorithm~1 (which cannot see the strata) and over a raw
combined decomposition (which cannot tell the strata from $\Var$ without
exactly this selector).

\subsection{Relation to the existing methods}

Phases I--II are a differential Thomas decomposition with the split set
generated lazily by the reduction trace --- a strategy, not a new
decomposition notion; the guarantees (disjointness, passivity of leaves) are
Thomas's \cite{BGLR2012}, and any correctness claim should lean on that
literature rather than re-prove it. Phase III$_\exists$ per leaf is the
bounded-prolongation route of \cite{bounded-prolongation-note} run per cell
(where its two known pathologies are mildest: on a cell the chain is
triangular, the inequations are explicit, and the closure order is the
cell's own). Phase III$_\forall$ is Algorithm~1 run relative to a cell, with
the saturation handled by radical membership. The prolong-to-order-$i$ /
algebraize / eliminate construction is classical (it is inside
\cite{GKO2009}, there with the $H$-saturation we deliberately omit); the
parametric splitting is comprehensive-Gr\"obner technology
\cite{weispfenning,gerdt-hashemi}. What this note adds is not machinery but
the \emph{quantifier discipline} --- the four-class table of
Definition~\ref{def:classes}, the two-pipelines theorem-in-outline of
\S\ref{sec:quantifier}, and Criterion~\ref{crit:main} resolving the defining%
-equation puzzle --- which is exactly the piece a correct implementation of
Brent's Phase~III sketch was missing.

%=====================================================================
\section{Open problems and next steps}
\label{sec:open}

\begin{enumerate}
\item \textbf{No a-priori closure order for $A\cup\{P\}$.} The ansatz bound
  $N_0(A)=2$ does not bound the combined system; the constants-cutting
  conditions $D^k\rho$ enter at order $2+k$ and the terminating $k$ is
  unknown for hydrogen (heuristically $\sim$4--6; unmeasured). The pointwise
  probe \code{prolong-probe.sage} is written to lower-bound it empirically
  but has not yet produced data \cite{existence-route-report}. An a-priori
  bound in general would bear on the effectivity gap of the Ritt problem
  \cite{GKO2008,GKO2009} --- do not expect one; prolong-until-stable is the
  algorithm.
\item \textbf{The parametric elimination cliff.} Direct III$_\exists$ over
  the eleven hydrogen constants as ring variables did not complete even over
  $\mathrm{GF}(32003)$ at order 2. The recommended path is per-cell staging
  over the 29 already-computed Thomas cells, with the cheap membership
  filter ($\Mem \subseteq \Ex$) applied first; syzygy-based consistency
  extraction is the principled single-ideal fallback
  \cite{existence-route-report}. Until one of these lands, the corrected
  algorithm's Phase III$_\exists$ is validated only on toy examples
  (\S\ref{sec:tests}).
\item \textbf{The ``used inequations'' completeness argument.} Phase II
  splits only on tracked divisors. The argument that this suffices ---
  untracked initials/separants do not affect the reduction identities, and
  the leaf's own solved-form initials are always used --- should be written
  out against B\"achler et al.'s simple-system conditions, in particular the
  squarefreeness splits (discriminants), which the tracking as described
  captures only when a gcd computation actually divides by them. Until then,
  treat Phases I--II as ``Thomas with a lazy strategy, conjecturally
  split-minimal,'' not as a proven refinement.
\item \textbf{Real witnesses.} III$_\exists$ certifies complex existence.
  For quantum-mechanical conclusions the real (and normalizability)
  questions are downstream filters on the resulting varieties, as in the
  paper's treatment; a systematic real-existential layer (real QE / CAD)
  is unexplored territory for this problem class.
\item \textbf{Verification against the recorded runs.} The concrete check
  that would close the loop: run III$_\exists$/III$_\forall$ per cell on the
  29 hydrogen cells and confirm (i) the membership cells' union reproduces
  \code{joca.sage}'s five primes, (ii) the partial strata include the
  $E=-1/2$ and $E=-1/8$ cells with their recorded primes, and (iii) nothing
  else appears. The classifier of \code{joca-rg-combined.sage} already
  implements the III$_\forall$ test and validates on the toy examples; the
  hydrogen combined RG run itself does not terminate, which is precisely why
  the per-cell route matters.
\item \textbf{$\mathrm{GF}(32003)$ vs $\mathbb{Q}$.} Any stratum destined for
  the paper needs a characteristic-zero confirmation of the modular
  computations.
\end{enumerate}

%=====================================================================
\begin{thebibliography}{99}

\bibitem{NewMethod} B.~Baccala, \emph{A New Method of Solving PDEs}, work in
progress, \code{\char`~/Papers/NewMethod/NewMethod.tex}. Algorithm~1 (core
solution: reduce and project), Theorem~3 (completeness), Algorithm~2
(comprehensive solution).

\bibitem{membership-note} \emph{Membership versus variety, and the
partial-solution strata}, companion note,
\code{membership-vs-variety-partial-strata.tex} (this directory).

\bibitem{bounded-prolongation-note} \emph{Bounded-Prolongation
Algebraization --- a proof attempt}, companion note,
\code{bounded-prolongation-theorem.md} (this directory).

\bibitem{method-comparison} \emph{NewMethod and the Parametric Decomposition
Landscape: a Comparison}, \code{method-comparison.md} (this directory),
2026-06-11.

\bibitem{litsearch} \emph{NewMethod parametric differential elimination ---
literature search}, project report
\code{\char`~/project/reports/newmethod-parametric-differential-elimination-litsearch.md},
2026-06-04.

\bibitem{existence-route-report} \emph{Bounded-Prolongation Existence Route:
2s Identification, Note Corrections, Implementation Findings}, project report
\code{\char`~/project/reports/bounded-prolongation-existence-route.md},
2026-07-07.

\bibitem{README} \code{README.md} (this directory): the \code{joca*.sage}
scripts, the RG performance findings, and the regularize-then-reduce
unfaithfulness record.

\bibitem{BGLR2012} T.~B\"achler, V.~Gerdt, M.~Lange-Hegermann, D.~Robertz,
\emph{Algorithmic Thomas decomposition of algebraic and differential
systems}, J.~Symbolic Comput.\ 47(10):1233--1266, 2012.

\bibitem{TDDS} V.~P.~Gerdt, M.~Lange-Hegermann, D.~Robertz, \emph{The MAPLE
package TDDS for computing Thomas decompositions of systems of nonlinear
PDEs}, Comput.\ Phys.\ Commun.\ 234:202--215, 2019.

\bibitem{LH2014} M.~Lange-Hegermann, \emph{The differential counting
polynomial}, Found.\ Comput.\ Math.\ (arXiv:1407.5838), 2014. Theorem~4.1.

\bibitem{blop} F.~Boulier, D.~Lazard, F.~Ollivier, M.~Petitot,
\emph{Computing representations for radicals of finitely generated
differential ideals}, AAECC 20(1):73--121, 2009.

\bibitem{fakouri} S.~Fakouri, S.~Rahmany, A.~Basiri, \emph{A new algorithm
for computing regular representations for radicals of parametric
differential ideals}, Cogent Mathematics \& Statistics 5(1):1507131, 2018.

\bibitem{GKO2008} O.~Golubitsky, M.~Kondratieva, M.~Moreno Maza,
A.~Ovchinnikov, \emph{A bound for the Rosenfeld--Gr\"obner algorithm},
J.~Symbolic Comput.\ 43(8):582--610, 2008.

\bibitem{GKO2009} O.~Golubitsky, M.~Kondratieva, A.~Ovchinnikov, \emph{On
the generalized Ritt problem as a computational problem}, J.~Math.\ Sci.\
163(5):515--522, 2009.

\bibitem{gerdt-hashemi} V.~Gerdt, A.~Hashemi, \emph{Comprehensive involutive
systems}, CASC 2012, LNCS 7442:123--138.

\bibitem{CTD2007} C.~Chen, O.~Golubitsky, F.~Lemaire, M.~Moreno Maza, W.~Pan,
\emph{Comprehensive triangular decomposition}, CASC 2007, LNCS 4770:73--101.

\bibitem{chen-morenomaza} C.~Chen, M.~Moreno Maza, \emph{Algorithms for
computing triangular decomposition of polynomial systems} (and the parametric
applications thereof), J.~Symbolic Comput.\ 47:610--642, 2012.

\bibitem{weispfenning} V.~Weispfenning, \emph{Comprehensive Gr\"obner bases},
J.~Symbolic Comput.\ 14:1--29, 1992.

\bibitem{montes-wibmer} A.~Montes, M.~Wibmer, \emph{Gr\"obner bases for
polynomial systems with parameters}, J.~Symbolic Comput.\ 45:1391--1425,
2010.

\bibitem{chaharbashloo} M.~S.~Chaharbashloo, A.~Basiri, S.~Rahmany,
S.~Zarrinkamar, \emph{An application of Gr\"obner basis in differential
equations of physics}, Z.~Naturforsch.\ A 68(10--11):646--650, 2013.

\bibitem{hubert} E.~Hubert, \emph{Notes on triangular sets and
triangulation-decomposition algorithms II: Differential systems}, LNCS
2630:40--87, 2003. (Ritt reduction, initials/separants, characterizable
components.)

\bibitem{seiler} W.~Seiler, \emph{Involution: The Formal Theory of
Differential Equations and its Applications in Computer Algebra}, Springer,
2010. (Formal integrability; coherence $\Leftrightarrow$ involution.)

\end{thebibliography}

\bigskip
\noindent\rule{0.3\textwidth}{0.4pt}\\
{\small\noindent Researched and written by Claude (Fable 5) on behalf of
Brent Baccala, task 463 of the project task runner; the worked examples were
verified by hand against the recorded outputs of \code{joca.sage},
\code{joca-rg.sage}, \code{joca-rg-combined.sage}, and the open-maple
differential Thomas run.\\
\code{/claude-fable-5}}

\end{document}
